\section{Optimization View: Local-Law-Compatible Summarizers}
\label{app:v7-optimization-projection}

The mergeable-summary discussion identifies the algebraic target: a tree
should build local states, merge those states, and apply the query only at the
root. This appendix spells out the corresponding optimization problem.
We choose a representable class of summarizers and train inside that
class while penalizing violations of the local laws. When a classical
or semantic target cannot be represented by a local-law-compatible
summarizer in that class, the remaining error is structural rather than
only an optimizer failure.

\subsection{The Objects Being Optimized}
\label{app:v7-optimization-objects}

Fix a tree \(T\) and an object space \(\mathcal X\). A node \(v\) in
the tree represents an object \(X_v\in\mathcal X\); leaves represent
raw blocks, and internal nodes represent the fixed composition of their
children. Let \(\Y\) be the task-value space with task distance
\(d_Y\), and let
\[
  \fstar:\mathcal X\to \Y
\]
be the oracle value we would like to preserve. In the experiments this
may be a classical sketch query, a Markov-count target, or an expert
score such as a manifesto dimension.

The learned tree has two maps. The summarizer
\[
  g_\theta:\mathcal I\to \mathcal Z
\]
is used both to summarize a raw leaf and to merge two child states by
calling it on their fixed composition \(z_L\oplus z_R\). Here
\(\mathcal I\) is the local input interface for \(g_\theta\): it
contains raw represented objects \(X\in\mathcal X\) and composed or
serialized state pairs \(z_L\oplus z_R\). The query
\[
  \widehat f_\phi:\mathcal X_{\mathrm{score}}\to\Y
\]
is evaluated on raw represented objects and on states through the same
fixed query interface. Concretely, either states already live in the
queryable domain, or a renderer
\(\rho:\mathcal Z\to\mathcal X_{\mathrm{score}}\) is applied before
the query is read; we suppress \(\rho\) and write
\(\widehat f_\phi(X)\) and \(\widehat f_\phi(z)\) directly. At the
root, this query turns the stored state into the final prediction.
The parameters
\(\theta\) may be weights of a neural operator, LoRA/adapter weights of
a transformer, or a prompt/program choice for an LLM-style operator. The
same notation covers all of these cases: the representable class is the
set
\[
  \mathcal C
  =
  \{\,g_\theta:\theta\in\Theta\,\}.
\]

For a document \(i\), the tree recursion is
\[
  z_{i,v}^{\theta}=g_\theta(X_{i,v})
    \quad\text{for leaves }v,
  \qquad
  z_{i,v}^{\theta}
    =
    g_\theta(z_{i,L}^{\theta}\oplus z_{i,R}^{\theta})
    \quad\text{for internal }v,
\]
and the root prediction is
\[
  \widehat y_i(\theta,\phi)
  =
  \widehat f_\phi(z_{i,\mathrm{root}}^{\theta}).
\]
The population target is the oracle answer
\(\fstar(X_{i,\mathrm{root}})\). The training problem is therefore to
choose \(g_\theta\in\mathcal C\) and \(\widehat f_\phi\) so that the
root prediction is accurate and the intermediate states behave like
mergeable summaries for that query.

For any query \(q:\mathcal X_{\mathrm{score}}\to\Y\), write
\[
  E_{\mathrm{root}}(g,q;T)
  =
  d_Y\!\left(q(z_{\mathrm{root}}^{g}),q(X_{\mathrm{root}})\right)
\]
for the root error of the tree summarizer \(g\) on tree \(T\), measured
under the query \(q\). When \(q=f^\star\), this is the true oracle root
error. When \(q=\widehat f\), it is the root error seen by the learned
teacher.

\subsection{Local Residuals}
\label{app:v7-local-residuals}

Let \(\ell_Y(y,y')\) be a nonnegative loss on task values, usually
\(d_Y(y,y')\) or \(d_Y(y,y')^2\). The three local laws become three
families of residuals.

\paragraph{C1: leaf sufficiency.}
For a leaf \(v\), the leaf state should preserve the answer the raw
block would receive:
\[
  r_{1,i,v}(\theta,\phi)
  =
  \ell_Y\!\left(
    \widehat f_\phi(g_\theta(X_{i,v})),
    \widehat f_\phi(X_{i,v})
  \right).
\]
If a gold oracle label is available at the leaf, the second argument may
be replaced by \(\fstar(X_{i,v})\). The displayed version is the
query-equivalence form used by the tree: it asks whether the learned
state is sufficient for the learned query.

\paragraph{C3: merge consistency.}
For an internal node \(v\) with children \(L,R\), the merged state
should preserve the answer for the parent object:
\[
  r_{3,i,v}(\theta,\phi)
  =
  \ell_Y\!\left(
    \widehat f_\phi(g_\theta(
      z_{i,L}^{\theta}\oplus z_{i,R}^{\theta})),
    \widehat f_\phi(X_{i,v})
  \right).
\]
This is the learned version of the mergeable-summary condition: valid
child states should merge into a valid parent state.

\paragraph{C2: re-summary stability.}
For any produced state \(z_{i,v}^{\theta}\), passing it back through
the summarizer should not change its answer-relevant content:
\[
  r_{2,i,v}(\theta,\phi)
  =
  \ell_Y\!\left(
    \widehat f_\phi(g_\theta(z_{i,v}^{\theta})),
    \widehat f_\phi(z_{i,v}^{\theta})
  \right).
\]
C2 is the stability law needed when summaries themselves become inputs
to later summarization calls.

Let \(\mathcal A_1,\mathcal A_3,\mathcal A_2\) be the audited or
training sets of C1, C3, and C2 units. If every unit is evaluated, the
empirical residuals are simple averages:
\[
  \widehat C_j(\theta,\phi)
  =
  \frac{1}{|\mathcal A_j|}
  \sum_{(i,v)\in\mathcal A_j}r_{j,i,v}(\theta,\phi),
  \qquad j\in\{1,2,3\}.
\]
Under a design-based audit, unit \((i,v)\) is sampled with known
probability \(\pi_{i,v}\). The corresponding inverse-probability
weighted form is
\[
  \widehat C_j^{\mathrm{IPW}}(\theta,\phi)
  =
  \frac{1}{|\mathcal U_j|}
  \sum_{(i,v)\in\mathcal U_j}
  \frac{A_{i,v}}{\pi_{i,v}}
  r_{j,i,v}(\theta,\phi),
\]
where \(A_{i,v}\in\{0,1\}\) records whether the unit was checked and
\(\mathcal U_j\) is the finite population of auditable units of type
\(j\). The optimization formulas below can use either the full
averages or their IPW versions; the difference is how the residuals are
estimated, not what they mean.

\subsection{The Empirical Objective}
\label{app:v7-training-objective}

The empirical objective is the loss used to fit the learned summarizer
and query. The summarizer \(g_\theta\) enters through the tree
recursion: changing \(\theta\) changes every produced state
\(z_{i,v}^{\theta}\). The query \(\widehat f_\phi\) enters when
we read the answer at the root, when we calibrate against available labels, and
when we measure C1/C3/C2 residuals through query equivalence.

The gold-standard loss measures whether the final tree prediction
matches the available document-level target:
\[
  \widehat L_{\mathrm{gold}}(\theta,\phi)
  =
  \frac{1}{n}
  \sum_{i=1}^n
  \ell_Y\!\left(
    \widehat f_\phi(z_{i,\mathrm{root}}^{\theta}),
    y_i
  \right).
\]
Here \(y_i\) is the available global label for document \(i\). In
settings with direct oracle access, \(y_i=\fstar(X_{i,\mathrm{root}})\).
In LLM and human-label settings, \(y_i\) is the observed global label
used to train a learned teacher.

The proxy-oracle gap loss trains the query against labeled objects,
usually full documents or raw concatenations during the first pass:
\[
  \widehat L_{\mathrm{gap}}(\phi)
  =
  \frac{1}{|\mathcal A_{\mathrm{gap}}|}
  \sum_{(a,y_a)\in\mathcal A_{\mathrm{gap}}}
  \ell_Y\!\left(
    \widehat f_\phi(a),
    y_a
  \right).
\]

The local-law part of the empirical loss is the weighted residual
\[
  \widehat E_{\mathrm{law},R}(\theta,\phi)
  =
  \rho_1\widehat C_1(\theta,\phi)
  +\rho_3\widehat C_3(\theta,\phi)
  +\rho_2(R-1)\widehat C_2(\theta,\phi),
\]
where the \(R-1\) factor records the number of possible re-summary
steps. The no-cost empirical objective is then
\begin{equation}
\label{eq:v7-learned-objective}
  \begin{aligned}
  \widehat L_{\mathrm{emp}}(\theta,\phi)
  =
  &\underbrace{\alpha_{\mathrm{gap}}\widehat L_{\mathrm{gap}}(\phi)}
    _{\text{proxy-oracle gap}}
  +\underbrace{\alpha_{\mathrm{gold}}\widehat L_{\mathrm{gold}}(\theta,\phi)}
    _{\text{root prediction}} \\
  &+\underbrace{\lambda\,\widehat E_{\mathrm{law},R}(\theta,\phi)}
    _{\text{local-law residuals}} .
  \end{aligned}
\end{equation}
The coefficients
\(\alpha_{\mathrm{gap}},\alpha_{\mathrm{gold}},\lambda,\rho_1,\rho_3,\rho_2\)
are training weights. They say how strongly the empirical procedure
cares about proxy-oracle agreement, distance to the gold standard, and
local-law violations. They are not error thresholds. The equivalent
notation with independent local-law weights
\(\lambda_1 C1+\lambda_3 C3+\lambda_2 C2\) is obtained by absorbing
\(\lambda\rho_j\) into the three coefficients.

\paragraph{Direct local-oracle regime.}
When local oracle labels are available, the local residuals may be
measured directly against \(\fstar\). For example, C1 can compare
\(\widehat f_\phi(g_\theta(b))\) or \(f(g_\theta(b))\) to
\(\fstar(b)\), and C3 can compare a merged parent state to the true
parent answer. This is the cleanest supervised version, and it matches
classical sketches whose local targets are exactly known.

\paragraph{Teacher-first LLM regime.}
When local oracle calls are unavailable, the training signal is
different. We first initialize the ladder with a simple summarizer:
the root state is the raw document or raw concatenation, so the first
query learner sees the same object that has the global label. Using
pairs \((X_i,y_i)\), we learn or select a teacher
\(\widehat f_\phi\) that predicts the global label on full documents.
After this first pass, \(\widehat f_\phi\) supplies the local
comparison used to train \(g_\theta\):
\[
  C1:\ \widehat f_\phi(g_\theta(b))\approx \widehat f_\phi(b),
  \qquad
  C3:\ \widehat f_\phi(g_\theta(z_L\oplus z_R))
      \approx \widehat f_\phi(X_L\oplus X_R),
  \qquad
  C2:\ \widehat f_\phi(g_\theta(z))\approx \widehat f_\phi(z).
\]
This does not claim that the teacher is the true oracle. It says that,
when local true labels are unavailable, the local laws are optimized in
the learned teacher's score space, and the final guarantee must add the
teacher's proxy-oracle gap.

\paragraph{From teacher residuals to oracle residuals.}
The conversion is just the triangle inequality. For any two objects or
states \(a,b\), write the teacher residual and true residual as
\[
  \widehat r(a,b)=d_Y(\widehat f(a),\widehat f(b)),
  \qquad
  r^\star(a,b)=d_Y(f^\star(a),f^\star(b)).
\]
If the teacher has pointwise oracle-prediction error
\[
  \eta(a)=d_Y(\widehat f(a),f^\star(a)),
\]
then
\[
  r^\star(a,b)
  \le
  \underbrace{\widehat r(a,b)}_{\text{teacher residual}}
  +
  \underbrace{\eta(a)+\eta(b)}_{\text{proxy-oracle gap slack}}.
\]
Thus a uniform proxy-oracle gap bound
\(\eta(\cdot)\le\varepsilon_{\mathrm{orc}}\) gives
\[
  r^\star(a,b)\le \widehat r(a,b)+2\varepsilon_{\mathrm{orc}}.
\]
Applied to C1, C3, and C2 separately, this says that the true-oracle
law residuals are controlled by the teacher-law residuals plus a
recovery charge for each local comparison. With the unweighted
composition below,
\[
  E_{\mathrm{law}}(g,f^\star;R)
  \le
  \underbrace{E_{\mathrm{law}}(g,\widehat f;R)}
    _{\text{teacher-law error}}
  +
  \underbrace{2(R+1)\varepsilon_{\mathrm{orc}}}
    _{\text{component proxy-oracle gap}},
\]
because the composed law error has one C1 term, one C3 term, and
\((R-1)\) C2 terms. This componentwise inequality is useful when we
want to interpret the local laws themselves in the true-oracle scale.
For final root prediction, the sharper route is to bound the root
distortion through \(\widehat f\) and then transfer that one root
comparison to \(f^\star\), which costs only
\(2\varepsilon_{\mathrm{orc}}\).

\subsection{Three Error Paths}
\label{app:v7-three-error-paths}

The same notation covers the three cases that arise in practice. Let
\(q:\mathcal X_{\mathrm{score}}\to\Y\) denote the query used to define
local-law residuals. In the direct-oracle case \(q=f^\star\). In the
teacher-first case \(q=\widehat f\). Let
\[
  \varepsilon_{\mathrm{score}}(q,f^\star)
\]
denote a uniform bound on the difference between \(q\) and the true
oracle on the objects and states compared by the tree. Thus
\(\varepsilon_{\mathrm{score}}(f^\star,f^\star)=0\), while the
teacher-first route has
\(\varepsilon_{\mathrm{score}}(\widehat f,f^\star)=\varepsilon_{\mathrm{orc}}\).
The first transfer is only calibration:
\[
  E_{\mathrm{root}}(g,f^\star;T)
  \le
  \underbrace{E_{\mathrm{root}}(g,q;T)}
    _{\text{checked root error}}
  +
  \underbrace{2\,\varepsilon_{\mathrm{score}}(q,f^\star)}
    _{\text{score approximation}}.
\]
If local laws also certify the learned-score root error by
\[
  E_{\mathrm{root}}(g,q;T)\le E_{\mathrm{law}}(g,q;R),
\]
then the corresponding true-oracle root error is bounded by
\[
  \underbrace{E_{\mathrm{law}}(g,q;R)}
    _{\text{local-law certificate}}
  +
  \underbrace{2\,\varepsilon_{\mathrm{score}}(q,f^\star)}
    _{\text{score approximation}}.
\]
This is the common certificate: root error for the checked score, plus
the price of using an estimated score instead of the true oracle.

The remaining distinction is whether \(f^\star\) itself is compatible
with mergeable state updates in the chosen class. In the online
non-mergeable case, training still chooses a summarizer \(g_\theta\) by
balancing root error against local-law error, now with the true oracle
\(f^\star\) available for both terms.
If \(f^\star\) has an exact mergeable witness in \(\mathcal C\), both
terms can in principle be driven to zero. If not, the learned tree is
the best tradeoff the chosen class can express; its remaining
\(E_{\mathrm{root}}(g_\theta,f^\star;T)\) and
\(E_{\mathrm{law}}(g_\theta,f^\star;R)\) are structural, not merely
optimizer noise.

\begin{table}[t]
\centering
\small
\setlength{\tabcolsep}{4pt}
\begin{tabular}{@{}p{0.23\linewidth}p{0.22\linewidth}p{0.25\linewidth}p{0.21\linewidth}@{}}
\toprule
Regime & Local-law query & Main error terms & Interpretation \\
\midrule
Only root labels are available &
\(q=\widehat f\), learned from global labels &
\(E_{\mathrm{law}}(g,\widehat f;R)+2\varepsilon_{\mathrm{orc}}\) &
We certify the tree through a calibrated teacher. Local checks are not
oracle calls. \\
Online oracle, mergeable target &
\(q=f^\star\) &
\(E_{\mathrm{law}}(g,f^\star;R)\), with zero possible when the class
contains the exact state law &
The local laws can hold exactly; the tree is estimating the oracle
itself. \\
Online oracle, non-mergeable target &
\(q=f^\star\), but no zero-gap mergeable \(g\) is available in
\(\mathcal C\) &
\(E_{\mathrm{root}}(g_\theta,f^\star;T)\) and
\(E_{\mathrm{law}}(g_\theta,f^\star;R)\), plus optimization, sampling,
and audit residuals &
The tree fits the best available root-accuracy/local-law tradeoff in
the chosen representation. \\
\bottomrule
\end{tabular}
\caption{The three error paths use the same local-law notation. The
difference is whether the local checks use the true oracle or a learned
teacher, and whether the true oracle lies in the mergeable summarizer
class being certified.}
\label{tab:v7-error-paths}
\end{table}

\subsection{Total Error, Certification, and Local-Law Pressure}
\label{app:v7-certified-objective}

The empirical objective is what training optimizes. The certificate is
stated separately in terms of audited residual bounds. Let
\[
  \varepsilon_{\mathrm{leaf}},\qquad
  \varepsilon_{\mathrm{merge}},\qquad
  \varepsilon_{\mathrm{idemp}}
\]
be audited upper bounds for the C1, C3, and C2 residual components.
For a tree run with \(R\) rounds, the local-law error is
\[
  E_{\mathrm{law}}(g,\widehat f;R)
  =
  \underbrace{\varepsilon_{\mathrm{leaf}}}_{\text{C1}}
  +
  \underbrace{\varepsilon_{\mathrm{merge}}}_{\text{C3}}
  +
  \underbrace{(R-1)\varepsilon_{\mathrm{idemp}}}_{\text{C2 reuse}}.
\]
If local laws are checked directly against the true oracle, the run is
certified at task accuracy \(\varepsilon\) when
\[
  E_{\mathrm{law}}(g,f^\star;R)\le \varepsilon.
\]
If the laws are checked through a learned teacher \(\widehat f\) whose
uniform proxy-oracle gap is \(\varepsilon_{\mathrm{orc}}\), the total certified
root-prediction error is
\[
  E_{\mathrm{tot}}(g,\widehat f;R)
  =
  \underbrace{E_{\mathrm{law}}(g,\widehat f;R)}
    _{\text{audited local laws}}
  +
  \underbrace{2\varepsilon_{\mathrm{orc}}}
    _{\text{root proxy-oracle gap}},
\]
and the acceptance condition is simply
\[
  E_{\mathrm{tot}}(g,\widehat f;R)\le \varepsilon.
\]

For set-restricted optimization arguments, it is useful to write an
idealized population objective. This is not the literal empirical
training loss; it is the abstraction saying that distance to the
gold-standard task value is balanced against local-law error. With
local-law shares
\(\rho_1,\rho_3,\rho_2\ge 0\), normalized so
\(\rho_1+\rho_3+\rho_2=1\), define
\[
  E_{\mathrm{law},\rho}(g,f;R)
  =
  \rho_1\varepsilon_{\mathrm{leaf}}
  +\rho_3\varepsilon_{\mathrm{merge}}
  +\rho_2(R-1)\varepsilon_{\mathrm{idemp}}.
\]
The balanced objective is
\begin{equation}
\label{eq:v7-oracle-projection-objective}
  \begin{aligned}
  J_{\lambda,\rho}(g;f)
  =
  &\underbrace{(1-\lambda)E_{\mathrm{root}}(g,f;T)}
    _{\text{root accuracy}} \\
  &+\underbrace{\lambda E_{\mathrm{law},\rho}(g,f;R)}
    _{\text{local-law pressure}} .
  \end{aligned}
\end{equation}
Here \(E_{\mathrm{root}}\) is the tree's root answer error. In the
direct regime, \(f=f^\star\). In the teacher-first regime,
\(f=\widehat f_\phi\) during training and the final true-oracle claim
adds the \(2\varepsilon_{\mathrm{orc}}\) proxy-oracle gap slack.
At \(\lambda=0\), this is pure gold-standard fitting. At
\(\lambda=1\), the objective minimizes only the local-law residual
penalty. The display explains the training pressure toward
local-law-compatible summarizers; it is not a claim that every
empirical optimizer minimizes exactly this displayed population
functional.

\subsection{Local-Law-Compatible Sets and Structural Error}
\label{app:v7-projection-sets}

The optimization view is clearest after separating the representable
class from the mergeable subset. Let \(\mathcal C\) be the class of
summarizers representable by the chosen neural operator, LLM prompt
family, or other learned operator. For a fixed tree \(T\) and query
\(f\), define the exact local-law set
\[
  \mathcal M_0(\mathcal C,T,f)
  =
  \{\,g\in\mathcal C:
      C1(g,f;T)=C2(g,f;T)=C3(g,f;T)=0\,\}.
\]
For approximate component bounds
\(\boldsymbol\varepsilon=(\varepsilon_{\mathrm{leaf}},
\varepsilon_{\mathrm{merge}},\varepsilon_{\mathrm{idemp}})\), define
\[
  \mathcal M_{\boldsymbol\varepsilon}(\mathcal C,T,f)
  =
  \{\,g\in\mathcal C:
      C1(g,f;T)\le\varepsilon_{\mathrm{leaf}},\
      C3(g,f;T)\le\varepsilon_{\mathrm{merge}},\
      C2(g,f;T)\le\varepsilon_{\mathrm{idemp}}\,\}.
\]
These are sets of functions, not necessarily linear spaces. The
word ``projection,'' when used below, means risk minimization over one
of these sets. If \(\mathcal R(g)\) is a target risk, the
set-restricted candidate is any minimizer of \(\mathcal R\) over
\(\mathcal M_{\boldsymbol\varepsilon}(\mathcal C,T,f)\).

\begin{proposition}[Set-restricted risk interpretation]
\label{prop:v7-projection-interpretation}
Set-restricted candidates minimize the chosen risk over the learnable
mergeable set. If that set contains a zero-risk representative and the
risk is nonnegative, every minimizer has zero residual gap. If every
member of the learnable mergeable set has positive risk, every
minimizer has a positive structural gap.
\end{proposition}

If the exact target is representable by a local-law-compatible
learned summarizer, set-restricted fitting can recover it in principle;
otherwise the best local-law-compatible learned summarizer still leaves
irreducible error. That error is structural: it comes from the chosen
representation class and local-law set, not from failing to optimize
hard enough.

\begin{proposition}[Local-law endpoint]
\label{prop:v7-local-law-endpoint}
Suppose the local-law penalty is faithful, meaning its zero set inside
\(\mathcal C\) is exactly the exact local-law set
\(\mathcal M_0(\mathcal C,T,f)\), and suppose the penalty is
nonnegative. If \(\lambda=1\) in the balanced objective and
\(\mathcal C\) contains at least one zero-penalty representative, then
any class-restricted minimizer lies in
\(\mathcal M_0(\mathcal C,T,f)\).
\end{proposition}

At the local-law endpoint, a faithful penalty forces minimizers onto
the exact mergeable set when that set is nonempty. For intermediate
\(0<\lambda<1\), the objective encourages this behavior but does not by
itself prove exact mergeability without additional convexity,
optimization, or margin assumptions.

\subsection{Proxy-Oracle Gap and Final Certification}
\label{app:v7-gap-certification}

The local laws are usually checked with the learned score
\(\widehat f_\phi\). The final scientific claim is about the true
oracle \(\fstar\). If the learned score has uniform proxy-oracle gap,
\[
  d_Y(\widehat f_\phi(a),\fstar(a))\le\varepsilon_{\mathrm{orc}}
  \qquad\text{for all relevant objects or states }a,
\]
then comparing through \(\widehat f_\phi\) costs at most
\(2\varepsilon_{\mathrm{orc}}\). Thus the operational certificate is
\[
  E_{\mathrm{law}}(g_\theta,\widehat f_\phi;R)+2\varepsilon_{\mathrm{orc}}
  \le
  \varepsilon
  \quad\Longrightarrow\quad
  \text{the tree is certified for }\fstar\text{ at accuracy }\varepsilon.
\]
The two terms have different meanings. \(E_{\mathrm{law}}\) is the
audited local-law error from C1, C3, and C2. The \(2\varepsilon_{\mathrm{orc}}\)
term is the price of using a calibrated learned score instead of the
oracle directly. Training weights such as \(\lambda\), \(\rho\), and
\(\alpha_{\mathrm{gap}}\) influence how the model is fit; they do not
replace the final inequality.

\subsection{How to Read Non-Mergeable Targets}
\label{app:v7-nonmergeable-targets}

The local-law-compatible set is most useful when the target is not a
known mergeable sketch. Suppose \(h^\star\) is an external behavior: a
non-mergeable library update rule, a one-way sketch being forced into
arbitrary tree merges, or a semantic oracle with no hand-designed
state. Training solves a constrained approximation problem:
\[
  \text{find the best }g_\theta\in\mathcal C
  \text{ that also lies near }
  \mathcal M_{\boldsymbol\varepsilon}(\mathcal C,T,\widehat f_\phi).
\]
If the gap after training and audit remains positive, the right
interpretation is that the chosen learned state class and local-law
constraints have found the closest mergeable approximation they can
represent. Increasing model capacity, changing the state format, or
changing the query may shrink that gap; lowering the optimizer
loss alone need not.

This is the optimization analogue of the mergeable-sketch story.
Classical sketches are the zero-gap case: the state, merge, and query
were designed so the local laws hold exactly. Learned C-Trees are the
learned case: the architecture supplies candidate summarizers, the
objective pushes them toward root accuracy and local-law compatibility,
and the audit decides whether the resulting tree is certified at the
target \(\varepsilon\).
